% Round 18. New-sectors branch. Candidate until the independent audit is bound.
\section{Fourfold hypersurfaces with a resolved hyperplane special fiber}
\label{r18hp:section}

This section constructs supplied strict semistable models with genuine
higher strata.  It uses only blowups with explicitly computed charts.
The resulting generic varieties need not have potential good reduction;
Section~\ref{r18hm:section} proves a stronger monodromy distinction when
there are at least six hyperplanes.  The model construction and the
associated edge-divisor generation calculation are separate arguments.

\subsection{The input and the ordered resolution}

Let $R$ be the valuation ring of a finite extension $K/\mathbf Q_p$,
let $\pi$ be a uniformizer, and let $k$ be its finite residue field.
Fix $d\ge3$ linear forms $\ell_1,\ldots,\ell_d$ on
$\mathbf P^5_R$ whose reductions are in geometric linear general
position.  Thus every intersection of at most five of their reduced
hyperplanes is a projective space of the expected dimension, and
any intersection of six is empty.  Write
\[
 P=\prod_{i=1}^d\ell_i,\qquad H_i=V(\bar\ell_i),\qquad
 H_I=\bigcap_{i\in I}H_i\simeq\mathbf P_k^{5-|I|}.
 \tag{HP18.1}\label{r18hp:arrangement}
\]
For a degree-$d$ form $Q$ over $R$, impose the following condition:
\[
 \begin{array}{l}
 V(\bar Q)\cap H_I\text{ is smooth of codimension one in }H_I
       \text{ if }1\le |I|\le4,\\
 V(\bar Q)\cap H_I=\varnothing\text{ if }|I|=5.
 \end{array}
 \tag{HP18.2}\label{r18hp:transverse}
\]
Empty strata are omitted.  In particular no positive-dimensional
$H_I$ is contained in $V(\bar Q)$.  Assume that the hypersurface
\[
                  X=V(P+\pi Q)\subset\mathbf P^5_K
 \tag{HP18.3}\label{r18hp:generic}
\]
is smooth and geometrically integral.

Start with $\mathcal M_0=\mathbf P^5_R$.  In this ambient scheme
blow up $H_1\subset(\mathcal M_0)_k$, then the strict transform
of $H_2$, and continue through the strict transform of $H_{d-1}$.
Denote the final ambient scheme by $\mathcal M$.  Let
$\mathcal X\subset\mathcal M$ be the closure of $X$; equivalently
it is the final strict transform of $V(P+\pi Q)$.
All blowup centers in the ambient schemes are the specified native
$k$-schemes, not centers chosen after an unnamed scalar extension.

\begin{theorem}[The native resolved model and every closed stratum]
\label{r18hp:model}
The scheme $\mathcal X$ is integral, regular, flat and projective
of relative dimension four over $R$.  Its special fiber is a
reduced strict normal-crossings divisor
\[
                         Y=\bigcup_{i=1}^d D_i.
 \tag{HP18.4}\label{r18hp:fiber}
\]
There are exactly $d$ irreducible components, indexed by the original
hyperplanes.  For every nonempty $I\subseteq\{1,\ldots,d\}$ with
$|I|\le5$, write $h=\max I$.  The closed stratum
$D_I=\bigcap_{i\in I}D_i$ is obtained from $H_I$ by blowing up,
in the order $j=d,d-1,\ldots,h+1$, the strict transforms of
\[
                    Z_{I,j}=H_I\cap H_j\cap V(\bar Q).
 \tag{HP18.5}\label{r18hp:centers}
\]
An empty center is skipped.  Every center which occurs is smooth;
its successive strict transforms are isomorphic to the original
$Z_{I,j}$.  In particular:
\[
 \begin{array}{c|c|c}
 |I|&H_I&\text{centers in }H_I\text{ and its successive blowups}\\ \hline
 1&\mathbf P^4&\text{smooth surfaces}\\
 2&\mathbf P^3&\text{smooth curves}\\
 3&\mathbf P^2&\text{smooth zero-dimensional schemes}\\
 4&\mathbf P^1&\text{none}\\
 5&\mathbf P^0&\text{none}.
 \end{array}
 \tag{HP18.6}\label{r18hp:types}
\]
Each $D_I$ is geometrically connected.  Intersections of six or
more components are empty.  Thus the dual complex is precisely
the $4$-skeleton of the $(d-1)$-simplex, with the convention that
it is the whole simplex when $d\le5$.
\end{theorem}

\begin{proof}
\emph{Step 1: native ambient centers.}
The center $H_i$ has codimension two in $\mathbf P^5_R$ and is
smooth over $k$.  Its strict transform after an earlier ambient
blowup is unchanged: the earlier center induces on $H_i$ the
ideal of the Cartier divisor $H_i\cap H_j$, and blowing up a
Cartier ideal is an isomorphism.  The same observation applies
successively.  In local coordinates the center of the next blowup
is given by $(t,x_i)$ below, so is a regular codimension-two
immersion in the already regular ambient scheme.  These equations
also verify the assertion at intersections of several earlier
centers.  Hence all ambient blowups are well defined and projective.

The original hypersurface is flat over $R$: its special equation
$P$ is nonzero, and its defining equation is primitive.  Its local
rings inject into their generic localizations.  The generic fiber
is integral, so the total hypersurface and its strict transforms
are integral.  A strict transform is computed by removing the
exceptional factor from the total-transform equation; the charts
below are the saturated charts, or equivalently closures of the
generic charts.  No component supported only on an exceptional
divisor is retained.

\emph{Step 2: the only local singularity to resolve.}
At a geometric point of the special fiber, let
$i_1<\cdots<i_s$ be exactly the hyperplanes passing through it.
At most five pass through a point.  Locally rename their equations
$x_1,\ldots,x_s$.  The other factors of $P$ are units.
If $\bar Q$ is nonzero at the point, the hypersurface is already
strict semistable, locally $x_1\cdots x_s=\pi$ after absorption of
a unit.  Each component ideal on this hypersurface is Cartier,
so the stated ambient blowups restrict to isomorphisms there.

Suppose now that $\bar Q=0$.  Condition
\eqref{r18hp:transverse} says that $x_1,\ldots,x_s$ and
$y=-Q/u$, where $u$ is the product of the unit factors, are part
of a relative etale coordinate system on the ambient smooth
scheme.  In particular $s\le4$ in this case.  The exact equation is
\[
                          x_1\cdots x_s=\pi y.
 \tag{HP18.7}\label{r18hp:binomial}
\]
The calculation is therefore in an actual etale chart, not merely
in a formal completion.  Extra smooth coordinates are suppressed.

Blow up $(\pi,x_1)$.  Its $\pi$-chart has
$x_1=\pi u_1$ and strict equation
\[
                         y=u_1x_2\cdots x_s.
\]
After eliminating $y$ this chart is regular and smooth over $R$.
Its other chart has $\pi=x_1t_1$ and strict equation
\[
                          x_2\cdots x_s=t_1y.
\]
The next ambient center is $(t_1,x_2)$ in this latter chart;
its strict transform is absent from the first chart.
Repeating gives the following complete collection of charts,
indexed by $m=1,\ldots,s$:
\[
 \begin{split}
  \pi&=x_1\cdots x_{m-1}t_{m-1},\\
  x_m&=t_{m-1}u_m,\\
  y&=u_mx_{m+1}\cdots x_s.
 \end{split}
 \tag{HP18.8}\label{r18hp:charts}
\]
For $m=1$, $t_0=\pi$ and the first line is tautological.
For $m=s$, set $u_s=y$: this is the final chart after the
first $s-1$ blowups, not an additional blowup.  Factors indexed
by global hyperplanes absent from the local point give identity
operations and are omitted.

The independent coordinates in chart $m$ are
$x_1,\ldots,x_{m-1},t_{m-1},u_m,x_{m+1},\ldots,x_s$
and the extra smooth coordinates, subject only to the first
line of \eqref{r18hp:charts}.  This is the standard regular
strict semistable chart.  Its special components are
$x_1=0,\ldots,x_{m-1}=0,t_{m-1}=0$, with multiplicity one.
Their generic points map birationally to the corresponding
original hyperplanes; the last divisor corresponds to $H_{i_m}$.
Thus the charts give regularity and reduced strict semistability,
and show that no additional vertical component occurs.
They also prove smoothness of every intersection.
Projectivity follows from the ambient construction, and
flatness follows from integrality and the nonconstant map to
$\operatorname{Spec}R$.

\emph{Step 3: identification of every stratum.}
Fix a set $I$ and let $h=\max I$.  In the local notation retain
the factors after the largest index in $I$.  The stratum $D_I$
is covered by precisely the charts in
\eqref{r18hp:charts} whose last displayed component has index
at least $h$.  On $D_I$ these charts, mapped to $H_I$, have the
following form: first the chart
\[
             y=u_h\prod_{j>h}x_j,
\]
and for each later factor $x_m$, the chart
\[
             x_m=t u_m,\qquad y=u_m\prod_{j>m}x_j.
 \tag{HP18.9}\label{r18hp:stratum-charts}
\]
The earlier $x_j$ not belonging to $I$ remain free coordinates.
These are exactly the charts obtained by blowing up on $H_I$
first $(y,x_d)$, then the strict-transform ideal $(y',x_{d-1})$,
and so on down to $(y^{\prime\cdots\prime},x_{h+1})$.
For example, the $x_d$-chart of the first blowup has $y=x_dy'$;
the other chart has $x_d=ty$ and is already the last chart of
\eqref{r18hp:stratum-charts}.  Induction gives every chart and
its transition functions.  The ideals are the global ideals of
\eqref{r18hp:centers}; the local descriptions therefore glue to
the asserted global iterated blowup, including over native
non-rational closed points.

By \eqref{r18hp:transverse}, each nonempty $Z_{I,j}$ is smooth
of dimension $3-|I|$.  Before blowing it up, a previously blown
center $Z_{I,j'}$, $j'>j$, induces on it the Cartier ideal of
$H_{j'}|_{Z_{I,j}}$.  This divisor is smooth whenever nonempty,
again by \eqref{r18hp:transverse}; the relevant local ideals
are $(x_j,y)$ and $(x_{j'},y)$.  Its strict transform is thus
isomorphic to $Z_{I,j}$, and the coordinate calculation shows
that it remains a regular codimension-two center in the
transformed $H_I$.  For zero-dimensional $Z_{I,j}$ these
intersections are empty, since $\bar Q$ avoids fivefold
hyperplane intersections.  For $|I|=4$ the proposed centers
are already empty for the same reason.  This proves
\eqref{r18hp:types} and the smoothness assertions.

Each iterated blowup is birational to the indicated projective
space, hence geometrically connected.  No intersection with
$|I|\ge6$ can occur, because its image would lie in the empty
$H_I$.  Every $|I|\le5$ stratum is nonempty by the blowup
description.  The last statement about the dual complex follows.
\end{proof}

\subsection{Numerical local exactness and ordinary actual Green forms}

The direct argument in HG18 computes the actual Picard groups of
the double strata and the actual triple divisor classes.  Their
exceptional-coordinate matrix is the negative identity.  For one
global relatively ample class $H$ it proves
$W^1=HW^0+\operatorname{im}\tau$, with the actual alternating
triple Gysin image contained in $\ker\sigma_2$.  The numerical
divisor-generation theorem DG18 then gives $U_2=0$.  This is the
primary argument; it requires no middle positivity on a component.

\begin{corollary}[Fourfold local exactness on the new models]
\label{r18hp:U2}
On the model of Theorem~\ref{r18hp:model}, with the notation and
signs of DQ17, one has
\[
 U_2=
 \frac{\ker(\rho_2:C^2\to W^2)\cap
                  \operatorname{im}(\sigma_2:W^1\to C^2)}
      {\operatorname{im}(\sigma_2\rho_1:C^1\to C^2)}=0,
 \qquad K_{\mathrm{comp},\ell}^2=\ker\rho_2.
 \tag{HP18.10}\label{r18hp:U2zero}
\]
Hence the ER17 refined quotient vanishes in codimension two.
More generally every specified geometrically homologically
trivial rational cycle on $X$ has exact-support ordinary actual
strong G in every codimension, over the original $K$ and for
all original compact actual tests.
\end{corollary}

\begin{proof}
Apply Theorem~\ref{r18hg:generation} in dimension four to
obtain $U_2=0$.  Its proof retains every triple and higher stratum
and uses the restrictions of one actual global ample class.

Every threefold double is an iterated blowup of $\mathbf P^3$
along smooth curves.  CD17 proves numerical-to-geometric-cycle-class
comparison for its curve classes, and this statement is retained
under the displayed smooth blowups over every finite extension of
$k$.  CL17 therefore gives the second equality of
\eqref{r18hp:U2zero}.  DQ17's exact sequence now makes the refined
model quotient zero, and ER17 proves obstruction zero for the
specified codimension-two cycle.  The other codimensions on a
fourfold are the already frozen divisor, one-cycle, and zero-cycle
cases.  The frozen original-field reconstruction gives the
ordinary actual Green form for the original specified cycle on
its exact support complement and its identity against every
compact actual test.  No analytic pushforward along the ambient
resolution appears in this argument.
\end{proof}

\paragraph{An independently checkable positivity route.}
The new signature theorem S18 also proves that $\mathbf P^4$
and its displayed iterated smooth blowups satisfy numerical
hard Lefschetz and primitive middle Hodge positivity for every
ample class.  The other strata have dimension at most three,
where the necessary numerical assertions follow from divisor
Hodge index.  Thus, conditional only on the separate proof and
review of S18, the frozen R6 numerical BGS theorem gives
$\ker A_a=\ker\rho_a+\operatorname{im}\sigma_a$ in every degree
on these same models.  This supplies a second route to $U_2=0$.
It is not a dependency of the direct HG18/DG18 proof above.

\begin{remark}[Extensions and the precise range of this theorem]
If $L/K$ is unramified, the same equations and the same model
satisfy the hypotheses over $L^\circ$.  Applying the theorem anew
to an arbitrary cycle on $X_L$ proves full G there; no claim that
such a cycle descends to $K$ is needed.  A further arbitrary
finite extension changes the coefficient of $Q$ to a unit times
a higher power of a uniformizer.  The present chart theorem does
not assert a semistable resolution for that new equation, and
this corollary alone does not establish the full-cycle theorem
on every ramified base change.  For classes actually pulled back
from $K$, the already frozen finite-trace and pullback interfaces
remain available with their full native degree.
\end{remark}

\paragraph{Nonemptiness and provenance.}
Over an algebraic closure, general hyperplanes satisfy
\eqref{r18hp:arrangement}, and the open set of degree-$d$ forms
satisfying \eqref{r18hp:transverse} is nonempty by Bertini on the
finitely many linear strata and avoidance of their zero-dimensional
members.  These data are defined over a finite extension of a
finite residue field, and over $\mathbf Q$ they can be chosen
rationally and then spread away from finitely many primes.
Smoothness of the generic member is an explicit hypothesis,
verified in the global construction below.
The use of ordered component blowups in hyperplane degenerations
is classical; for example Gross--Siebert discuss it in
\emph{Affine manifolds, log structures, and mirror symmetry}.
No priority claim is made for the resolution procedure.
The contribution here is its explicit native fourfold stratum
calculation and activation of the refined numerical criterion
by the direct edge-divisor generation calculation in HG18.
